\subsection{借助项来定义函数}

C54. 设 T、A 为两个项，且 x、y 为不同的字母。假设 x 不出现在 A 中，且 y 既不出现在 T 中也不出现在 A 中。设 R 为关系“ \(x \in A\) 且 \(y = T\) ”。关系 R 关于字母 x 和 y 有一个图 F。该图是函数图；其第一投影为 A，第二投影为当 \(x \in A\) 时所有形如 T 的对象所成的集合（§1，第6号）。对于每一个 \(x \in A\) 都有 \(F(x) = T\) .

设 B 为当 \(x \in A\) 时所有形如 T 的对象所成的集合。于是

\[
\boldsymbol {R} \Longrightarrow ((\boldsymbol {x}, \boldsymbol {y}) \in \boldsymbol {A} \times \boldsymbol {B});
\]

由于由 \(A \times B\) 既不包含x也不包含y，R关于字母x和y有一个图F（第1号）。显然，该关系

\[
(x, y) \in F \text {   且   } (x, y ^ {\prime}) \in F
\]

蕴含 \(y = y'\) ，因此 F 是一个函数图。其余断言显然。

如果C是一个集合，它包含由形如T的对象组成的集合B， \(x \in A\) （其中y不出现在C中），则函数 \((F, A, C)\) 也记作记号 \(x \to T\) ( \(x \in A\) , \(T \in C\) ）；相应的形式数学中的汇编既不包含x也不包含y，并且不依赖于满足上述条件的字母y的选择。当上下文足够明确时，我们将允许使用记号 \(x \to T(x \in A)\) , \((T)_{x \in A}\) ，或 \(x \to T\) ，有时简写为 T 或 \((T)\) 。*因此我们可以称“函数 \(x^{3}\) ''，如果上下文明确表明我们指的是该映射 \(x \to x^{3}\) 将复数集合映射到自身的函数。*

\minorheading{示例}

\begin{enumerate}
\def\labelenumi{(\arabic{enumi})}
\item
  如果 \(f\) 是从 \(\mathbf{A}\) 到 \(\mathbf{B}\) 的映射，则函数 \(f\) 等于函数 \(x \to f(x)\) ( \(x \in \mathbf{A}\) , \(f(x) \in \mathbf{B}\) ），后者简写为 \(x \to f(x)\) ，也可写作 \((f_x)_{x \in \mathbf{A}}\) （后一记号尤其用于“元素族”而不是“函数”这一说法）。
\item
  让 \(\mathbf{G}\) 是有序对的一个集合。函数
\end{enumerate}

\[
z \rightarrow \operatorname{pr} _ {1} z (z \in \mathbf {G}, \operatorname{pr} _ {1} z \in \operatorname{pr} _ {1} \mathbf {G})
\]

并且

\[
z \rightarrow \mathrm{pr} _ {2} z \quad (z \in \mathrm{G}, \mathrm{pr} _ {2} z \in \mathrm{pr} _ {2} \mathrm{G})
\]

分别称为G上的第一坐标函数和第二坐标函数；当没有混淆风险时，它们被记为 \(pr_{1}\) 并且 \(pr_{2}\) 。

